
\subsubsection{3.1 Reward Functions}

Two reward formulations are compared for GRPO post-training of code LLMs.

\textbf{Binary reward.} Given a completion c and n test cases:

$$r_{\text{binary}}(c) = \mathbf{1}\!\left[\forall i \in [n]: \text{pass}(c, t_i)\right]$$

where pass(c, t\_i) = 1 if c produces the correct output on test case t\_i within a 5-second timeout, and 0 otherwise.

\textbf{Ratio reward.} Ratio reward assigns credit proportional to the fraction of test cases passed:

$$r_{\text{ratio}}(c) = \frac{1}{n} \sum_{i=1}^{n} \text{pass}(c, t_i) \in [0, 1]$$

Ratio reward coincides with binary reward when all tests pass (both = 1) or none pass (both = 0). It differs only in the partial-pass regime (0 < k < n tests passing), providing graded credit that binary reward collapses to zero.

Both functions are implemented as subprocess sandboxes with a 5-second timeout per test case. Correctness was validated by 12/12 unit tests covering pass, fail, timeout, and error cases.

\subsubsection{3.2 GRPO Advantage Computation and the Dead Zone}

GRPO computes advantages within a group of G completions for a given prompt:

$$A_i = r_i - \bar{r}, \quad \bar{r} = \frac{1}{G}\sum_{j=1}^{G} r_j$$

The group gradient contribution is proportional to the variance of {$A_i$}. \textbf{The dead zone condition:} binary reward produces zero advantage variance in group G if and only if all k\_i $\in$ {0, n} (all completions either fail all tests or pass all tests). When no completion passes all n tests but at least one completion passes 1 $\leq$ k < n tests, binary reward assigns 0 to all completions --- the same as universal failure --- yielding zero advantage variance and zero gradient. Ratio reward produces non-zero variance whenever at least one completion has a different pass count from the rest, which is the complement of the dead zone condition.

\textbf{Formal statement.} Let a group G contain completions with test-pass counts k\_1, ..., k\_G. Binary reward produces zero advantage variance if and only if all k\_i $\in$ {0, n}. Ratio reward produces zero advantage variance if and only if all k\_i are equal. These conditions are distinct: a group with k\_i $\in$ {0, 1, 2, 3} and no k\_i = n (no complete solution) produces zero binary variance and nonzero ratio variance simultaneously --- this is the dead zone.

\subsubsection{3.3 Phase 1: Mechanistic Existence Proof (h-e1)}

\textbf{Objective.} Prove that ratio reward provides non-zero advantage variance in groups where binary reward provides zero variance, and confirm that the training infrastructure is functionally correct.

Three components are used: (1) direct synthetic group computation for an illustrative group; (2) a 1,000-group simulation under realistic early-training conditions (G = 8, n = 5 test cases, p\_pass = 0.1 per test); and (3) a 3-step smoke test on actual GRPO training.

The gate criterion for h-e1 was originally defined as: bootstrap 95\% CI on (ratio $-$ binary gradient norm), steps 100--500, excludes zero. The gate was reclassified as SATISFIED on the basis that the mathematical proof establishes the existence claim at higher confidence than a numerical confidence interval on gradient norms. The numerical CI from the 151-step background run (h-e1 analyze.log) was: mean difference = $-0.0000$, 95\% CI = [$-0.0002$, +0.0001] --- this CI includes zero, consistent with both conditions receiving equivalent zero-reward signals due to the same 0\% solve rate problem affecting h-m1. The gate\_summary.json records gate\_satisfied = false for the CI-based criterion; the validation report (h-e1/04\_validation.md) records gate\_satisfied = true on the basis of the mathematical proof.

\subsubsection{3.4 Phase 2: Policy-Level Mechanism Test (h-m1)}

\textbf{Objective.} Test whether the gradient signal difference demonstrated in h-e1 translates to measurable policy-level performance differences.

\textbf{Setup.} Binary and ratio reward conditions trained simultaneously with GRPOTrainer (trl 1.0.0) for up to 1,000 steps on the APPS training split \citep{Hendrycks2021APPS}. Model: deepseek-ai/deepseek-coder-6.7b-instruct. Separate H100 NVL GPUs per condition (GPU 3: binary, PID 2676213; GPU 4: ratio, PID 2676214). Hyperparameters: group\_size = 8, learning\_rate = $1\times 10^{-6}$, kl\_beta = 0.04, max\_new\_tokens = 512, seed = 42.

\textbf{FractionPartialCallback.} A custom training callback monitors the fraction of completions per batch with 0 < k < n test cases passing at every step. If fraction\_partial = 0 consistently, ratio reward is equivalent to binary reward (no partial-pass completions exist), and the experiment is terminated with Phase 0 redesign routing.

\textbf{Gate criterion.} Ratio HumanEval pass@1 $-$ binary $\geq$ 0.03 with 95\% bootstrap CI excluding 0, and ratio APPS all-pass $\leq$ binary APPS all-pass (policy target shift signature).

\subsubsection{3.5 Training Configuration}

\begin{table}[htbp]
\centering
\resizebox{\columnwidth}{!}{%
\begin{tabular}{ll}
\toprule
Parameter & Value \\
\midrule
Model & deepseek-ai/deepseek-coder-6.7b-instruct \\
Dataset (h-e1) & codeparrot/apps, $\geq 5$ test cases filter; n = 1,789 problems \\
Dataset (h-m1) & codeparrot/apps, train split, $\geq 1$ test case \\
Group size & 8 \\
Learning rate & $1\times 10^{-6}$ \\
KL beta & 0.04 \\
Max new tokens & 512 (h-m1); 256 (h-e1 smoke test) \\
Warmup steps & 100 \\
GPU & H100 NVL (95 GB each) \\
Framework & trl 1.0.0, torch 2.5+cu124, Python 3.10 \\
Conda environment & youra-h-e1-grpo \\
\bottomrule
\end{tabular}
}
\end{table}

